% @author Lluis Gesa (lluis.gesa@uab.cat)
%
% Enginyeria Del Software
% -----------------------------------------------------------------------------------



% definim quin estil general volem per el document
\documentclass[11pt]{article}

% indiquem quins paquets amb funcionalitat extra volem fer servir
\usepackage{a4wide, amsmath, tabularx, colortbl, fancyhdr, graphicx, lastpage, caption}
\usepackage[utf8]{inputenc}

% Indiquem que comencem a definir el document i els seus continguts.
\begin{document}

% Definim valors i propietats, creem macros que ens ajudaran a
% agilitzar coses repetides:

\newcommand{\NomProjecte}{Perdidos en la ignorancia}
\newcommand{\NomDocument}{Document de Joc}
\newcommand{\IdentificadorDocument}{GG-SL-N-GD}
\newcommand{\VersioDocument}{v1.0}

% Definim una macro que cada cop que insertem \requeriment, creara una
% taula amb el contingut passat per paràmetre
\newcommand{\requeriment}[6]{
\begin{center}
	\begin{tabularx}{\linewidth}{|p{2.5cm}|X|}
	\hline
	\cellcolor[gray]{.8} \textbf{ID} & \cellcolor[gray]{.9} #1\\
	\hline
	\cellcolor[gray]{.8} \textbf{Title} & #2 \\
	\hline
	\cellcolor[gray]{.8} \textbf{Description} & #3\\
	\hline
	\cellcolor[gray]{.8} \textbf{Priority} & #4 \\
	\hline
	\cellcolor[gray]{.8} \textbf{Verification} & #5 \\
	\hline
	\cellcolor[gray]{.8} \textbf{Parents} & #6\\
	\hline
	\end{tabularx}
\end{center}
}

% Fent servir propietas del paquet fancyhdr per crear una capçalera
% que es mostrarà cada pàgina
\pagestyle{fancy}
\fancyhead[LO,LE]{\includegraphics[width = 3.5cm]{./imatges/uab-logo.png}}
\fancyhead[CO,CE]{\bf \NomProjecte \\ \vspace{\baselineskip}
  \NomDocument }
\fancyhead[RO,RE]{
{\bf Ref : \IdentificadorDocument} \\
Versió: \VersioDocument\\
Data: \today \\
Pag. \thepage ~of \pageref{LastPage}
}
\fancyfoot[LO,LE,CO,CE,RO,RE]{}
\textheight=19cm


%
% Creem la pàgina inicial :

% Insertem el nom del document en una taula centrada feng servir una
% font gran \huge

\begin{center}
\begin{tabular*}{15cm}{p{14.5cm}}
\\
\\
\\
\rowcolor[gray]{.9}\begin{center}\huge{\bf{\NomDocument}}\end{center}\\
\\
\end{tabular*}
\end{center}


%
% Insertem una taula que contindrà les abreviacions que hi hagi al
% nostre document
%
\small{%
\begin{center}
\begin{tabular}{|p{2cm}|p{4.75cm}|p{2cm}|p{4.75cm}|}
\hline
\multicolumn{4}{|>{\columncolor[gray]{.9}[2mm][2mm]}c|}{{\bf Abreviatures}} \\
\hline
UAB & Universitat Autonoma de Barcelona & ES & Enginyeria del Software \\
\hline
& & & \\
\hline
\end{tabular}
\end{center}
}

%
% Insertem una taula amb el historial de canvis del document
%
\small{%
\begin{center}
\begin{tabular}{|p{2cm}|p{2cm}|p{7.5cm}|p{2cm}|}
\hline
\multicolumn{4}{|>{\columncolor[gray]{.9}[2mm][2mm]}c|}{{\bf Historial de Revisions}} \\
\hline
{\bf Version} & {\bf Date} & {\bf Comments} & {\bf Autor} \\
\hline
1.0 & 08-Abril-2019 & Sprint 1 &
G43-2-02 \\
\hline
\end{tabular}
\end{center}
}




% Insertem en una nova pàgina, els indexs : de continguts, de figures
% i de taules, les comandes \tableofcontents ,  \listoffigures,
% \listoftables són gestionades per els paquets fets servir, i per
% defecte escriuen els titols de la secció en angles, per canviar-los,
% hem de 'renovar les comandes' :

\renewcommand{\contentsname}{Índex de continguts}
\renewcommand{\listfigurename}{Índex de figures}
\renewcommand{\listtablename}{Índex de taules}

% Ara li diem al LaTeX que inserti la informació
\newpage
\tableofcontents
\listoffigures
\listoftables

% Insertem les seccions que volem
\newpage

\section{Introducció}\label{sec:intro}

\begin{flushleft}
En este documento se pretende detallar de forma minuciosa el comportamiento deseado de todos los casos de uso del sistema. De esta forma, será mucho más fácil saber que debe hacer cada módulo cuando se quiera dividir el conjunto de tareas a programar.\\


Además, esto ayuda a anticiparse a posibles necesidades y problemas, lo que agiliza el proyecto y ayuda a determinar con mayor exactitud las herramientas a utilizar y tiempo a invertir.\\


En este mismo documento se proporcionará los diagramas de casos de uso, para detectar de forma sencilla los actores que intervienen y sus respectivas acciones.\\
\end{flushleft}

\section{Descripció}

\subsection{extres}

\section{Personatges}
En el juego solo habrá un tipo de personaje, que será el jugador. Dado la clase de juego que es no es necesaria la inclusión de variantes.
\section{Històries}

Se han recogido todas las necesidades para llevar a cabo el desarrollo.

-	El usuario debe conectarse al sistema mediante un usuario y una contraseña.

-	Para pasar de nivel los usuarios deben completar los objetivos del nivel.

-	El juego tiene un ranking que se podrá consultar desde internet.

-	En las habitaciones se pueden encontrar objetivos o “bonus” que ayudarán al usuario a pasarse los posteriores niveles o acceder a los siguientes.

-	En algunos niveles aparecen pruebas contrarreloj que obligan al usuario a pasarlo en tiempo o serán eliminados cuando finalice el tiempo establecido por el nivel.

-	Si el jugador es eliminado en una prueba contrareloj, debe volver a empezar la partida desde el principio.

-	El juego tiene los roles de administrador y de jugador.

-	El administrador puede crear y editar nuevos mapas.

-	El administrador puede modificar la contraseña del usuario.

-	El administrador puede eliminar la cuenta del usuario.

-	El administrador puede modificar el ranking de los usuarios.

-	El usuario puede modificar o eliminar sus datos de la cuenta.

-	El sistema se encarga de gestionar todo lo relacionado con el juego (rankings, cuentas, niveles...).

-	El sistema es el encargado de generar los mapas aleatoriamente y sus propiedades (tipo de salas o habitaciones, la dificultad del nivel, tipo de nivel: normal o contrarreloj).

-	Hay cierta cantidad de mapas de cada dificultad, que son seleccionados por el sistema de forma aleatoria en función del nivel en el que te encuentras.

-	El sistema gestiona los niveles de menor a mayor dificultad.

-	El administrador puede ver el estado de los usuarios que se encuentran en linea (en partida...).

-	El administrador puede crear cuentas nuevas.

-	El jugador puede continuar la partida por donde la dejó.

-	El sistema se encarga de guardar el estado de la partida del usuario.

-	El jugador puede empezar una partida nueva.

-	El jugador puede borrar una partida existente.

-	El jugador puede consultar a los bonus obtenidos anteriormente, que serán mostrados por pantalla y puede seleccionar entre utilizar uno de ellos o no usar ninguno.

-	En la pantalla se mostrará al jugador si se encuentra en una pantalla contrarreloj y el tiempo del cual dispone para pasarse el nivel.

-	En cada turno el jugador puede consultar todas las acciones que tiene disponibles en la propia pantalla (menú de movimientos).

-	El jugador puede guardar la partida manualmente desde el menú de la partida.

-	El jugador tiene en el menú de la partida la opción para salir del juego.

-	El jugador tiene en el menú de la partida la opción para volver al menú principal.

-	El jugador tiene en el menú principal la opción para modificar sus ajustes.

-	El jugador tiene en el menú principal la opción para ver los rankings.

-	El jugador tiene en el menú principal la opción para empezar una partida nueva.

-	El jugador tiene en el menú principal la opción para gestionar sus partidas existentes.

-         El jugador puede acceder a los movimientos que ha realizado previamente durante la partida.

-	El jugador finalizará la partida si completa todos los niveles existentes en la partida.



% Finalitzem el document.
\end{document}
